% Irodov Part III Electrodynamics §3.1, original 1--12.
% Source PDF pp.101--102 (printed pp.92--93), answers PDF pp.304--305.
% k=1/(4*pi*epsilon_0). Part III RESTARTS independent numbering.
\electroproblem{1}{92}
求两电子、两质子之间静电力与万有引力之比。
若两相同粒子的静电力和万有引力大小相等，
其电荷质量比 $|q|/m$ 应为多少？
\exsolution
库仑力与引力大小为
$F_e=e^2/(4\pi\varepsilon_0r^2)$、
$F_g=Gm^2/r^2$。
故比值与距离无关：
\[
\frac{F_e}{F_g}=\frac{e^2}{4\pi\varepsilon_0 Gm^2}.
\]
代入 $e=1.602\times10^{-19}\,\mathrm C$、
$m_e=9.109\times10^{-31}\,\mathrm{kg}$、
$m_p=1.673\times10^{-27}\,\mathrm{kg}$，
得到
\[
\boxed{(F_e/F_g)_{ee}\simeq4.2\times10^{42}},\quad
\boxed{(F_e/F_g)_{pp}\simeq1.2\times10^{36}}.
\]
一般相同粒子两力相等要求
$q^2/(4\pi\varepsilon_0r^2)=Gm^2/r^2$，
故
\[
\boxed{|q|/m=\sqrt{4\pi\varepsilon_0G}
\simeq8.6\times10^{-11}\,\mathrm{C/kg}}.
\]
\excheck \textbf{底本勘误：}
书末把质子比值印作 $10^{26}$，
与所列基本常数相差10个数量级；
应为约 $10^{36}$，两种力均按真空中的点粒子处理。

\electroproblem{2}{92}
两个铜球，质量各为 $1.0\,\mathrm g$，
球心相距 $1.0\,\mathrm m$。
若每球的电子总电荷与核内正电荷之差
为核电荷的 $1\%$，
求它们之间的电相互作用力数量级。
\exsolution
铜原子序数 $Z=29$、
摩尔质量 $M_{\rm Cu}\simeq63.5\,\mathrm{g/mol}$。
单球的原子数为
$N=(m/M_{\rm Cu})N_A$，
其核正电荷总量为 $ZNe$。
净电荷模
\[
|Q|=0.01\,ZNe
=0.01\cdot29\cdot\frac{1.0}{63.5}
(6.022\times10^{23})(1.602\times10^{-19})
\simeq4.4\times10^2\,\mathrm C.
\]
用库仑定律
\[
\boxed{F=\frac{Q^2}{4\pi\varepsilon_0r^2}
\simeq1.8\times10^{15}\,\mathrm N
\approx2\times10^{15}\,\mathrm N}.
\]
\excheck 这是概念性的巨大电荷差假设；
实际金属球不可能在普通环境下
稳定承载数百库仑的同号电荷。

\electroproblem{3}{92--93}
两相同小球，各质量 $m$、电荷 $q(t)$，
由长为 $\ell$ 的等长绝缘线悬挂于同一点。
两球距离 $x\ll\ell$。
若两球靠近速率满足
$-\dot x=v=a/x$，求电荷泄漏速率 $\dot q$。
\exsolution
在小角度准静态平衡下
$\tan\theta\simeq x/(2\ell)$，
两球间静电排斥力
$F_e=q^2/(4\pi\varepsilon_0x^2)$，
竖直平衡给线张力约为 $mg$，
因此
\[
\frac{q^2}{4\pi\varepsilon_0x^2}
=mg\frac{x}{2\ell}
\quad\Longrightarrow\quad
q=\sqrt{\frac{2\pi\varepsilon_0mg}{\ell}}\,
x^{3/2}.
\]
对时间求导并用 $\dot x=-a/x$：
\[
\boxed{\dot q=
-\frac{3a}{2}
\sqrt{\frac{2\pi\varepsilon_0mg}{\ell x}}}.
\]
泄漏量率为
$-\dot q=(3a/2)\sqrt{2\pi\varepsilon_0mg/(\ell x)}$。
\excheck 若原书抽取答案将 $x^{-1/2}$ 漏掉，
会导致量纲不等于 $\mathrm{C/s}$；
上式由逐步微分和单位检查确认。

\electroproblem{4}{93}
两个正点电荷 $q_1,q_2$，
位置矢量分别为 $\bm r_1,\bm r_2$。
求应放置在何处、大小为何的第三个负点电荷
$q_3$，使三个电荷各自受合力为零。
\exsolution
要使两正电荷的排斥力被吸引力抵消，
$q_3<0$ 必须处于两者连线之间。
设两正电荷间距 $L$，第三电荷距 $q_1$ 为 $x$。
第三电荷自身零受力要求
$q_1/x^2=q_2/(L-x)^2$，
得
\[
x=\frac{\sqrt{q_1}}{\sqrt{q_1}+\sqrt{q_2}}L,
\quad
\boxed{\bm r_3=
\frac{\sqrt{q_2}\,\bm r_1+\sqrt{q_1}\,\bm r_2}
{\sqrt{q_1}+\sqrt{q_2}}}.
\]
对 $q_1$ 的力平衡：
$q_1q_2/L^2=q_1|q_3|/x^2$，
故
\[
\boxed{q_3=-\frac{q_1q_2}
{(\sqrt{q_1}+\sqrt{q_2})^2}}.
\]
\excheck 对称特例 $q_1=q_2=q$，
$q_3=-q/4$ 且位于中点，
与库仑力的距离平方关系相符。

\electroproblem{5}{93}
半径 $r$ 的细导线圆环总电荷为 $q$。
若在圆环中心加入点电荷 $q_0$，
求导线的\emph{拉伸张力增量} $\Delta\mathcal T$。
\exsolution
环上每单位弧长的电荷
$\lambda=q/(2\pi r)$。
取中心角为 $\mathrm d\varphi$ 的微小圆环元，
其电荷
$\mathrm dq=\lambda r\,\mathrm d\varphi
=q\,\mathrm d\varphi/(2\pi)$。
中心电荷对该弧元施加径向力
$\mathrm dF=(1/(4\pi\varepsilon_0))
q_0\mathrm dq/r^2$。
两端张力增量的径向合力
$2\Delta\mathcal T\sin(\mathrm d\varphi/2)
\simeq\Delta\mathcal T\,\mathrm d\varphi$。
比较两式：
\[
\boxed{\Delta\mathcal T=
\frac{q q_0}{8\pi^2\varepsilon_0r^2}}.
\]
\excheck 相同符号电荷导致张力增加，
不同符号则为负增量；
这里求的是新增张力而非环自身静电排斥
原有的总张力。

\electroproblem{6}{93}
真空中有正点电荷 $q=50\,\mathrm{\mu C}$，
位于二维位置
$\bm r_0=(2\bm i+3\bm j)\,\mathrm m$。
求在
$\bm r=(8\bm i-5\bm j)\,\mathrm m$
处的电场矢量及模。
\exsolution
位移矢量
$\bm R=\bm r-\bm r_0
=(6\bm i-8\bm j)\,\mathrm m$，
其模 $R=10\,\mathrm m$。
点电荷电场
\[
\bm E=\frac{q}{4\pi\varepsilon_0}
\frac{\bm R}{R^3}.
\]
代入
$1/(4\pi\varepsilon_0)\simeq8.988\times10^9$
得
\[
\boxed{\bm E\simeq
(2.7\bm i-3.6\bm j)\,\mathrm{kV/m}},
\quad
\boxed{|\bm E|\simeq4.5\,\mathrm{kV/m}}.
\]
\excheck 负 $y$ 分量反映测量点位于电荷下方；
模满足 $\sqrt{2.7^2+3.6^2}=4.5$。

\electroproblem{7}{93}
图3.1所示的正方形两上顶点各放置 $+q$，
两下顶点各放置 $-q$，
正方形两条对角线长均为 $2\ell$。
求从正方形中心沿其法线距离 $x$ 处的
电场强度模。
\exsolution
取正方形位于 $z=0$，
上、下顶点纵坐标为
$\pm\ell/\sqrt2$，
各顶点到轴上观察点
$(0,0,x)$ 的距离均为
$\sqrt{\ell^2+x^2}$。
库仑场叠加时，
水平方向分量相互抵消，
轴向分量因总电荷为零而抵消，
只余从正电荷向负电荷方向的平面内分量。
单个顶点对该分量的贡献
为 $k q(\ell/\sqrt2)/(\ell^2+x^2)^{3/2}$，
四项相加：
\[
\boxed{E(x)=
\frac{q\ell}{\sqrt2\,\pi\varepsilon_0
(\ell^2+x^2)^{3/2}}}.
\]
\excheck 该场\emph{平行正方形所在平面}，
而不是垂直于平面；
$x\gg\ell$ 时呈偶极型 $x^{-3}$ 衰减。

\electroproblem{8}{93}
一条薄半圆弧半径 $R=20\,\mathrm{cm}$，
沿弧均匀带总电荷 $q=0.70\,\mathrm{nC}$。
求其曲率中心处电场强度模。
\exsolution
弧长 $\pi R$，线密度
$\lambda=q/(\pi R)$。
沿半圆以角度 $\varphi$ 参数化，
弧元
$\mathrm dq=\lambda R\,\mathrm d\varphi$，
电场大小微元为
$\mathrm dE=k\,\mathrm dq/R^2$。
平行弧直径的场分量两两抵消，
对称轴分量求和：
\[
E=\frac{k\lambda}{R}
\int_{-\pi/2}^{\pi/2}\cos\varphi\,\mathrm d\varphi
=\frac{2k\lambda}R
=\boxed{\frac{q}{2\pi^2\varepsilon_0R^2}}.
\]
代入得
\[
\boxed{E\simeq1.0\times10^2\,\mathrm{V/m}
=0.10\,\mathrm{kV/m}}.
\]
\excheck 完整均匀圆环的中心电场为零，
但半圆环沿其对称轴的场不为零。

\electroproblem{9}{93}
半径为 $r$ 的均匀带电细圆环总电荷为 $q$。
求圆环轴线上距中心 $\ell$ 处的电场模，
讨论 $\ell\gg r$，求最大值及其位置，
并说明 $E(\ell)$ 的大致图像。
\exsolution
每个环元到轴上点距离
$s=\sqrt{r^2+\ell^2}$，
径向分量由对称性相消，
仅保留轴向分量：
\[
E(\ell)=\frac{1}{4\pi\varepsilon_0}
\frac{|q\ell|}{(r^2+\ell^2)^{3/2}}.
\]
在 $\ell\gg r$ 时，
$E\simeq|q|/(4\pi\varepsilon_0\ell^2)$。
对 $\ell>0$ 求导：
\[
E'(\ell)\propto
\frac{r^2-2\ell^2}{(r^2+\ell^2)^{5/2}},
\]
故最大值出现在 $\ell=r/\sqrt2$，
\[
\boxed{E_{\max}=
\frac{|q|}{6\sqrt3\,\pi\varepsilon_0r^2}}.
\]
图像从 $E(0)=0$ 起先上升，
到 $\ell=r/\sqrt2$ 达峰值后按
$\ell^{-2}$ 衰减。
\excheck 环轴电场在中心恰为零；
若把环误当成点电荷，会在 $\ell=0$
得到错误的发散。

\electroproblem{10}{93}
一点电荷 $+q$ 处于均匀带电细圆环中心，
环半径为 $R$、总电荷为 $-q$。
求在轴线上 $x\gg R$ 处的电场模。
\exsolution
在 $x>0$ 处，
中心点电荷的轴向场
$E_+=kq/x^2$，
负圆环的轴向场
$E_-=-kq x/(x^2+R^2)^{3/2}$。
所以
\[
E_x=\frac{kq}{x^2}
\left[1-\left(1+\frac{R^2}{x^2}\right)^{-3/2}\right].
\]
当 $R/x\ll1$，
展开
$(1+u)^{-3/2}=1-\frac32u+O(u^2)$，
得到
\[
\boxed{E\simeq
\frac{3qR^2}{8\pi\varepsilon_0x^4}}.
\]
\excheck \textbf{原书书末系数争议：}
本版简答印
$3qR^2/(4\pi\varepsilon_0x^4)$，
是上述逐项叠加与二项式展开结果的两倍。
按照题面点电荷 $+q$ 与圆环总电荷 $-q$
的条件，本稿取 $3/(8\pi)$ 系数，
待独立对照原版勘误或其他版本。

\electroproblem{11}{93}
一个半径 $R$、总电荷为 $q$ 的均匀带电细圆环，
与一根从圆环中心沿轴向延伸至远方的
均匀带电细长直线同轴；
长直线从圆环中心开始，
线电荷密度为 $\lambda$。
求两者间的相互作用力。
\exsolution
从长直线起点起沿轴向取坐标 $z\ge0$。
线元电荷为 $\lambda\,\mathrm dz$，
到圆环任意环元的距离为
$\sqrt{R^2+z^2}$。
线元对圆环元的\emph{径向}场分量为
\[
\mathrm dE_R=
\frac1{4\pi\varepsilon_0}
\frac{\lambda R\,\mathrm dz}{(R^2+z^2)^{3/2}}.
\]
径向分量对整个均匀环对称相消；
更方便的办法是求\emph{轴向作用力}：
圆环轴线上电场为
$E_z(z)=kqz/(R^2+z^2)^{3/2}$，
长直线上受力为
\[
F=\lambda\int_0^\infty E_z(z)\,\mathrm dz
=kq\lambda
\int_0^\infty\frac{z\,\mathrm dz}{(R^2+z^2)^{3/2}}
=\boxed{\frac{q\lambda}{4\pi\varepsilon_0R}}.
\]
\excheck 作用力方向沿公共轴，
符号由 $q\lambda$ 决定；
使用牛顿第三定律，圆环受力大小相同方向相反。

\electroproblem{12}{93}
半径为 $R$ 的绝缘细圆环线电荷密度
$\lambda(\varphi)=\lambda_0\cos\varphi$。
求电场模：(a) 圆心；
(b) 环轴线上距圆心 $x$ 处；
并求 $x\gg R$ 的渐近表达式。
\exsolution
取 $\varphi=0$ 沿平面内 $+X$ 轴。
电荷元
$\mathrm dq=\lambda_0R\cos\varphi\,\mathrm d\varphi$，
$\int_0^{2\pi}\mathrm dq=0$。
观察点在轴线上时，场的轴向分量与总电荷成正比，
故为零；平面内 $X$ 分量为
\[
E_X=-k\frac{R}{(x^2+R^2)^{3/2}}
\int_0^{2\pi}\cos\varphi\,\mathrm dq
=-k\frac{\pi\lambda_0R^2}{(x^2+R^2)^{3/2}}.
\]
所以
\[
\boxed{E(0)=\frac{|\lambda_0|}{4\varepsilon_0R}},
\qquad
\boxed{E(x)=
\frac{|\lambda_0|R^2}{4\varepsilon_0(x^2+R^2)^{3/2}}}.
\]
定义环的偶极矩
$p=\int(R\cos\varphi)\,\mathrm dq
=\pi\lambda_0R^2$。
在 $x\gg R$ 时
\[
\boxed{E\simeq\frac{|p|}{4\pi\varepsilon_0x^3}}.
\]
\excheck 由于线密度正负交替，
该环\emph{净电荷为零}；
轴上场方向在环平面内，
表现为偶极场而非点电荷的 $x^{-2}$ 场。
